\documentclass[11pt]{article}
\usepackage[margin=1in]{geometry}
\usepackage{booktabs}
\usepackage{graphicx}
\usepackage{hyperref}
\usepackage{enumitem}
\setlength{\parindent}{0pt}
\setlength{\parskip}{0.6em}

\title{Prompt-Profile Binary Retrain: \texttt{h256 d2} On Balanced \texttt{majority\_s\_0.5}}
\author{}
\date{2026-04-04 06:18 UTC}

\begin{document}
\maketitle

\section*{Executive Summary}
\begin{itemize}[leftmargin=1.5em]
\item This note covers a new binary-only retrain on the saved April \texttt{majority\_s\_0.5} data. It is not a new target-choice run and it does not overwrite the earlier locked-pair report.
\item Exact runtime:
  \begin{itemize}[leftmargin=1.5em]
  \item pilot \texttt{2106}: \texttt{2026-04-04 05:54:44 UTC} to \texttt{2026-04-04 05:56:55 UTC}
  \item corrected all-dataset rerun \texttt{2107}: \texttt{2026-04-04 05:57:19 UTC} to \texttt{2026-04-04 06:04:55 UTC}
  \end{itemize}
\item Fixed object: same saved prompt-prefill activations, same binary target \texttt{majority\_s\_0.5}, same train-balanced / test-natural split contract, same optimizer seeds \texttt{0,1,2}, same \texttt{epochs=15}, \texttt{lr=1e-4}, \texttt{weight\_decay=0.05}, \texttt{dropout=0.1}.
\item The only intentional model change was probe capacity: old default \texttt{hidden\_dim=128}, \texttt{depth=1}; retrain \texttt{hidden\_dim=256}, \texttt{depth=2}.
\item Main read on held-out \texttt{PR-AUC}: \texttt{h256 d2} is not a uniform win. Ensemble improves on \texttt{GPQA}, \texttt{MATH-500}, and \texttt{MMLU-Pro}, but drops on \texttt{AIME} and \texttt{LiveCodeBench}. Last-layer improves on \texttt{GPQA}, \texttt{MMLU-Pro}, and \texttt{LiveCodeBench}, but drops on \texttt{AIME} and \texttt{MATH-500}.
\item Best current global ranking surface from this rerun: \texttt{ensemble h256 d2}. Mean test \texttt{PR-AUC = 0.492}, versus \texttt{0.474} for default ensemble, \texttt{0.430} for \texttt{last\_layer h256 d2}, and \texttt{0.350} for the prompt-length baseline.
\item Limitation: this rerun tuned probe capacity only. It did not sweep which layers to keep. \texttt{ensemble} still means one MLP per layer across all \texttt{28} saved layers, and \texttt{last\_layer} still means final layer only.
\end{itemize}

\section*{Exact Question}
If we keep the saved April balanced binary data fixed, does moving from the inherited default probe family to a larger \texttt{h256 d2} family materially improve the binary \texttt{majority\_s\_0.5} screen?

\section*{Run Contract}
\begin{itemize}[leftmargin=1.5em]
\item Source root: \texttt{/data/scratch/murphy/outputs/cot-loop-detection/full\_train\_locked\_pair\_20260404}
\item Retrain root: \texttt{/data/scratch/murphy/outputs/cot-loop-detection/full\_train\_locked\_pair\_20260404\_binary\_h256d2}
\item Binary target: \texttt{majority\_s\_0.5}
\item Train balance: downsample to \texttt{50/50}
\item Test balance: natural
\item Views:
  \begin{itemize}[leftmargin=1.5em]
  \item \texttt{ensemble}: one MLP per layer, aggregated by \texttt{vote\_fraction}
  \item \texttt{last\_layer}: final-layer-only MLP
  \end{itemize}
\item Important bug note:
  \begin{itemize}[leftmargin=1.5em]
  \item \texttt{2106} only covered \texttt{LiveCodeBench}, because repeated \texttt{--dataset} flags collapsed to the last dataset only
  \item commit \texttt{913021c} fixed that bug, and \texttt{2107} is the real five-dataset rerun
  \end{itemize}
\end{itemize}

\section*{What Changed Versus The April Default}
\begin{center}
\small
\setlength{\tabcolsep}{4pt}
\begin{tabular}{@{}lll@{}}
\toprule
Setting & Default binary run & Retrain \\
\midrule
Hidden dim & \texttt{128} & \texttt{256} \\
MLP depth & \texttt{1} & \texttt{2} \\
Dropout & \texttt{0.1} & \texttt{0.1} \\
Epochs & \texttt{15} & \texttt{15} \\
Learning rate & \texttt{1e-4} & \texttt{1e-4} \\
Weight decay & \texttt{0.05} & \texttt{0.05} \\
Layer set & all \texttt{28} layers for ensemble; final layer for \texttt{last\_layer} & unchanged \\
Data / split & saved April balanced binary data & unchanged \\
\bottomrule
\end{tabular}
\end{center}

So this is a capacity comparison, not a new data object and not a layer-subset sweep.

\section*{Ranking Results: Test \texttt{PR-AUC}}
\begin{center}
\small
\setlength{\tabcolsep}{4pt}
\resizebox{\textwidth}{!}{
\begin{tabular}{@{}lrrrrrr@{}}
\toprule
Dataset & Test prevalence & Prompt length & Ensemble \texttt{h128 d1} & Ensemble \texttt{h256 d2} & Last-layer \texttt{h128 d1} & Last-layer \texttt{h256 d2} \\
\midrule
GPQA & 0.0500 & 0.0657 & 0.4198 & 0.5025 & 0.2298 & 0.3176 \\
AIME & 0.5833 & 0.8976 & 0.9218 & 0.9090 & 0.9043 & 0.8479 \\
MATH-500 & 0.0400 & 0.1002 & 0.1354 & 0.1596 & 0.1506 & 0.1322 \\
MMLU-Pro & 0.0125 & 0.1104 & 0.1837 & 0.2023 & 0.0559 & 0.0959 \\
LiveCodeBench & 0.3375 & 0.5760 & 0.7111 & 0.6865 & 0.7116 & 0.7575 \\
\bottomrule
\end{tabular}
}
\end{center}

\section*{Cross-Dataset Mean \texttt{PR-AUC}}
\begin{center}
\small
\begin{tabular}{@{}lr@{}}
\toprule
Surface & Mean test \texttt{PR-AUC} \\
\midrule
Prompt length baseline & 0.3500 \\
Default ensemble \texttt{h128 d1} & 0.4744 \\
Default \texttt{last\_layer h128 d1} & 0.4104 \\
Retrain ensemble \texttt{h256 d2} & 0.4920 \\
Retrain \texttt{last\_layer h256 d2} & 0.4302 \\
\bottomrule
\end{tabular}
\end{center}

\section*{Retrain Threshold Behavior}
These threshold metrics are secondary here because the train split is balanced while the test split is natural. On the rare-positive datasets, high recall / low precision is expected.

\begin{center}
\small
\setlength{\tabcolsep}{4pt}
\resizebox{\textwidth}{!}{
\begin{tabular}{@{}llrrrrl@{}}
\toprule
Dataset & View & Precision & Recall & Positive-F1 & Macro-F1 & Selected epochs \\
\midrule
GPQA & ensemble & 0.1538 & 1.0000 & 0.2667 & 0.5487 & 4;7;5 \\
GPQA & last\_layer & 0.1065 & 1.0000 & 0.1862 & 0.3375 & 6;15;6 \\
AIME & ensemble & 1.0000 & 0.4762 & 0.6330 & 0.6840 & 4;9;4 \\
AIME & last\_layer & 0.8778 & 0.6667 & 0.7564 & 0.7494 & 12;6;15 \\
MATH-500 & ensemble & 0.0761 & 1.0000 & 0.1414 & 0.4008 & 12;8;8 \\
MATH-500 & last\_layer & 0.0759 & 1.0000 & 0.1407 & 0.3842 & 15;7;11 \\
MMLU-Pro & ensemble & 0.0794 & 1.0000 & 0.1470 & 0.5336 & 12;1;7 \\
MMLU-Pro & last\_layer & 0.0330 & 0.5000 & 0.0611 & 0.4920 & 15;15;3 \\
LiveCodeBench & ensemble & 0.6766 & 0.5494 & 0.5895 & 0.7057 & 3;3;14 \\
LiveCodeBench & last\_layer & 0.6014 & 0.7901 & 0.6828 & 0.7396 & 15;15;15 \\
\bottomrule
\end{tabular}
}
\end{center}

\section*{Interpretation}
\begin{itemize}[leftmargin=1.5em]
\item Proven now:
  \begin{itemize}[leftmargin=1.5em]
  \item \texttt{h256 d2} raises the best cross-dataset mean \texttt{PR-AUC} for both views
  \item \texttt{ensemble h256 d2} is the strongest overall ranking surface in this rerun
  \item within the retrained family, ensemble beats \texttt{last\_layer} on \texttt{4 / 5} datasets; \texttt{LiveCodeBench} is the exception
  \end{itemize}
\item Also proven:
  \begin{itemize}[leftmargin=1.5em]
  \item the effect is not monotone by dataset
  \item \texttt{AIME} is worse under the larger family for both views on \texttt{PR-AUC}
  \item \texttt{LiveCodeBench} splits the other way: larger \texttt{last\_layer} is the best surface there, while larger ensemble is worse than the default ensemble
  \end{itemize}
\item Important threshold caveat:
  \begin{itemize}[leftmargin=1.5em]
  \item rare-positive datasets (\texttt{GPQA}, \texttt{MATH-500}, \texttt{MMLU-Pro}) naturally push the retrained classifiers toward high recall and weak precision under the train-fit threshold
  \item that is why \texttt{PR-AUC} remains the primary comparison object for this rerun
  \end{itemize}
\item What this rerun did not answer:
  \begin{itemize}[leftmargin=1.5em]
  \item whether we should drop early layers, keep only a subset of layers, or tune the layer set per dataset
  \item this pass only changed MLP capacity, not the ensemble layer membership
  \end{itemize}
\end{itemize}

\section*{Recommendation}
\begin{itemize}[leftmargin=1.5em]
\item For the current balanced binary object, keep \texttt{PR-AUC} as the main decision metric.
\item If one single global surface is needed today, use \texttt{ensemble h256 d2}.
\item If we are willing to split by view or by dataset:
  \begin{itemize}[leftmargin=1.5em]
  \item keep the ensemble family for \texttt{GPQA}, \texttt{AIME}, \texttt{MATH-500}, and \texttt{MMLU-Pro}
  \item keep \texttt{last\_layer h256 d2} as the strongest current \texttt{LiveCodeBench} binary surface
  \end{itemize}
\item The next honest tuning step, if we want more than this, is not another vague ``bigger probe'' change. It is a small layer-subset / view sweep on the same balanced binary data.
\end{itemize}

\section*{Artifact Bundle}
\begin{itemize}[leftmargin=1.5em]
\item Metrics CSV: \texttt{outputs/prompt\_profile\_binary\_retrain\_h256d2\_20260404/binary\_retrain\_metrics.csv}
\item Summary JSON: \texttt{outputs/prompt\_profile\_binary\_retrain\_h256d2\_20260404/binary\_retrain\_summary.json}
\item Slurm logs:
  \begin{itemize}[leftmargin=1.5em]
  \item \texttt{outputs/prompt\_profile\_binary\_retrain\_h256d2\_20260404/logs/prompt-profile-binary-retrain-2106.out}
  \item \texttt{outputs/prompt\_profile\_binary\_retrain\_h256d2\_20260404/logs/prompt-profile-binary-retrain-2107.out}
  \end{itemize}
\end{itemize}

\end{document}
